\documentclass[11pt,a4paper]{article}
\usepackage[utf8]{inputenc}
\usepackage[T1]{fontenc}
\usepackage[french]{babel}
\usepackage{lmodern}
\usepackage[margin=2.1cm]{geometry}
\usepackage{booktabs,array,tabularx}
\usepackage{xcolor}
\usepackage{fancyvrb}
\usepackage{fvextra}
\usepackage{tikz}
\usetikzlibrary{arrows.meta,positioning,fit}
\usepackage[hidelinks]{hyperref}
\usepackage{textcomp}
\usepackage{newunicodechar}
\newunicodechar{→}{\ensuremath{\rightarrow}}
\newunicodechar{₂}{\textsubscript{2}}
\usepackage{amsmath,amssymb}
\newunicodechar{✓}{\ensuremath{\checkmark}}
\newunicodechar{✗}{\ensuremath{\times}}
\newunicodechar{Ω}{\ensuremath{\Omega}}
\newunicodechar{−}{-}


\setlength{\parskip}{0.4em}
\setlength{\parindent}{0pt}
\fvset{fontsize=\small,frame=single,framesep=2mm,rulecolor=\color{black!30},breaklines=true,breakanywhere=true}
\usepackage{xurl}
\DeclareUrlCommand\vp{\urlstyle{tt}}  % chemins et identifiants (coupure possible)
\newcommand{\cmd}[1]{\texttt{\detokenize{#1}}}  % commandes ASCII avec espaces (pas d'accents : detokenize)
\newcommand{\aremplacer}[1]{\textcolor{red!60!black}{\texttt{#1}}}
\newcommand{\attention}[1]{\par\noindent\fbox{\parbox{0.96\linewidth}{\textbf{Attention.} #1}}\par}
\renewcommand{\arraystretch}{1.15}
\newcolumntype{L}{>{\raggedright\arraybackslash}X}

\title{Intégrer VetTwins au matériel d'une salle de simulation\\[0.3em]
\large Relais, moniteur, auscultation, injection, machine d'anesthésie, pouls, boîtier du formateur}
\author{Projet VetTwin}
\date{20 novembre 2026}

\begin{document}
\sloppy
\maketitle

\begin{center}
\fbox{\parbox{0.92\linewidth}{\textbf{Usage pédagogique uniquement.} VetTwins est un patient virtuel canin destiné à
l'enseignement ; beaucoup de paramètres du module d'anesthésie sont des placeholders (niveau de preuve D). Aucun
résultat ne doit servir à décider d'une dose ou d'un traitement réel. Les boîtiers décrits ici sont des accessoires
d'entraînement à très basse tension (USB 5~V) : aucun n'est un dispositif médical.}}
\end{center}

\tableofcontents

\section{Principe}
\subsection{Qui calcule quoi}
La salle reprend l'onglet \emph{Anesthésie} de l'application web (\vp{web/anesthesia.js}). Trois règles
structurent l'installation :
\begin{enumerate}
  \item \textbf{Un seul moteur.} Seule la \emph{régie} (le navigateur du formateur, page \vp{index.html?room=1})
  exécute la session d'anesthésie. Les autres postes affichent ou jouent ce qu'elle diffuse ; le lecteur
  d'auscultation est la seule exception, il synthétise le son lui-même à partir du patient et de la FC reçus.
  \item \textbf{Les boîtiers mesurent, la régie traduit.} Un boîtier envoie ce qu'il constate dans la salle (seringue
  de propofol à \(10\)~mg/mL, \(6\)~mL poussés ; cadran du vaporisateur à \(2\)~\% ; oxymètre posé). La régie convertit
  en événement du moteur (mg/kg avec le poids du patient) et \emph{renvoie au boîtier le verdict du moteur} : action
  acceptée, ou refusée avec son motif (« intubation impossible : déglutition présente »). Aucun boîtier ne connaît
  la physiologie ni le poids.
  \item \textbf{Un relais entre les deux.} Un petit serveur (\vp{salle/relay.py}, Python, bibliothèque standard)
  sert l'application, fait transiter les messages et garde le dernier état pour les boîtiers de sortie (pouls).
  Il ne calcule rien.
\end{enumerate}

\begin{figure}[h]
\centering
\resizebox{\linewidth}{!}{%
\begin{tikzpicture}[font=\small,
  b/.style={draw, rounded corners, align=center, minimum height=9mm, inner sep=4pt},
  a/.style={-{Stealth[length=2mm]}, thick}, l/.style={font=\footnotesize, fill=white, inner sep=1.5pt}]
\node[b, fill=blue!8] (relay) at (0,0) {\textbf{Relais} \texttt{salle/relay.py}\\ Mac mini, Raspberry Pi ou portable\\ \texttt{http://IP:8765}};
\node[b, fill=green!10] (regie) at (0,4.2) {\textbf{Régie} (formateur)\\ \texttt{index.html?room=1}\\ moteur WebAssembly};
\node[b, fill=red!6] (mon) at (7.2,1.3) {\textbf{Moniteur de salle}\\ \texttt{monitor.html}};
\node[b, fill=red!6] (ausc) at (7.2,-1.5) {\textbf{Lecteur d'auscultation}\\ \texttt{ausc.html} \(\to\) haut-parleur\\ du thorax};
\node[b, fill=orange!10] (in) at (-7.4,1.3) {\textbf{Boîtiers d'entrée} (ESP32)\\ capteurs posés, stéthoscope,\\ pupitre RFID, machine,\\ boîtier du formateur};
\node[b, fill=orange!10] (out) at (-7.4,-1.5) {\textbf{Boîtier de sortie} (ESP32)\\ pouls fémoral};
\draw[a, <->] (regie) -- node[l, left=2pt, align=right]{WebSocket \texttt{/ws}\\ entrées et verdicts \(\downarrow\uparrow\), état diffusé} (relay);
\draw[a] (relay.east) -- node[l, above, sloped]{état, tracés} (mon.west);
\draw[a] (relay.east) -- node[l, below, sloped]{patient, site, FC} (ausc.west);
\draw[a, <->] (in.east) -- node[l, above, sloped]{\texttt{POST /api/device}} (relay.west);
\draw[a] (relay.west) -- node[l, below, sloped]{\texttt{GET /api/pulse}} (out.east);
\end{tikzpicture}}
\caption{Architecture de la salle. Toutes les liaisons restent sur le réseau local de la salle.}
\end{figure}

\subsection{Contenu du dépôt}
\begin{tabularx}{\linewidth}{>{\raggedright\arraybackslash}p{0.38\linewidth}L}
\toprule
\textbf{Fichier} & \textbf{Rôle} \\
\midrule
\vp{salle/relay.py} & relais : fichiers de \vp{web/}, WebSocket \vp{/ws}, API des boîtiers \\
\vp{salle/vettwins-relay.service}, \vp{salle/fr.vettwins.relay.plist} & lancement au démarrage (systemd, launchd) \\
\vp{web/room.js} & liaison des pages au relais (reconnexion automatique) \\
\vp{web/anesthesia.js} (fin du fichier) & régie : traduction des entrées des boîtiers, diffusion de l'état \\
\vp{web/monitor.html}, \vp{monitor.js}, \vp{monitor-view.js} & moniteur de salle ; moniteur défilant partagé avec la régie \\
\vp{web/ausc.html}, \vp{ausc.js} & lecteur d'auscultation \\
\vp{salle/firmware/libraries/VetTwinRelay} & bibliothèque Arduino commune (Wi-Fi, envoi, lecture du pouls) \\
\texttt{salle/firmware/<boîtier>/}\allowbreak\texttt{<boîtier>.ino} & micrologiciel de chaque boîtier (six croquis) \\
\vp{tests/check_relay.py}, \vp{check_firmware.py}, \vp{room_e2e.mjs} & essais (section~\ref{sec:essais}) \\
\bottomrule
\end{tabularx}

\subsection{Ce qui a été vérifié, et ce qui ne l'a pas été}
\begin{itemize}
  \item \textbf{Vérifié} (20/11/2026) : le relais (\vp{tests/check_relay.py}, 39 contrôles : aller-retour boîtier
  \(\to\) régie \(\to\) verdict, refus, délai, clé, sortie du dossier web refusée) ; la chaîne complète dans Chromium
  (\vp{tests/room_e2e.mjs} : régie, moniteur, lecteur d'auscultation, boîtiers simulés par HTTP, refus d'une
  intubation trop précoce renvoyé tel quel au boîtier) ; les six croquis compilés sur ordinateur en C++11 contre
  des en-têtes factices, et la bibliothèque \vp{VetTwinRelay} exécutée sur des réponses réelles du relais
  (\vp{tests/check_firmware.py}).
  \item \textbf{Non vérifié} : la compilation pour ESP32 avec la chaîne Arduino réelle, le téléversement et le
  câblage. Aucun boîtier n'a été monté. Les essais de la section~\ref{sec:essais} sont à faire à la première mise
  en route ; signalez les écarts.
\end{itemize}

\section{Réseau et relais}\label{sec:relais}
\subsection{Matériel}
\begin{itemize}
  \item \textbf{Serveur} : le Mac mini déjà décrit dans \vp{docs/DEPLOIEMENT_MACMINI.pdf}, un Raspberry Pi 4 ou 5,
  ou un portable. Le relais ne calcule rien : toute machine qui exécute Python~3.8 convient.
  \item \textbf{Point d'accès Wi-Fi dédié à la salle}, en 2,4~GHz (l'ESP32 n'a pas de 5~GHz), protégé par WPA2
  personnel. Les ESP32 ne savent pas se connecter simplement aux réseaux d'établissement à identifiant
  (WPA2-Entreprise, eduroam) : d'où ce point d'accès. Il n'a pas besoin d'accès à Internet.
  \item Une \textbf{adresse fixe} pour le serveur (réservation DHCP dans le point d'accès), par exemple
  \aremplacer{192.168.50.10}. Les écrans et le poste de régie peuvent être en Wi-Fi ou en Ethernet sur le même
  réseau.
\end{itemize}

\subsection{Installation}
\begin{Verbatim}
cd ~/srv/VetTwin/vetsim-kit
sh web/build.sh                       # moteur WebAssembly et modules (voir DEPLOIEMENT_MACMINI)
python3 -I salle/relay.py --web web --host 0.0.0.0 --port 8765
\end{Verbatim}
\cmd{--host 0.0.0.0} écoute sur toutes les interfaces : le relais n'est à lancer ainsi que sur le réseau isolé de la
salle. Option \cmd{--key SECRET} (ou variable \vp{VETTWIN_KEY}) : clé partagée exigée par l'API et le WebSocket ; les
pages la reçoivent par \vp{&key=SECRET}, les boîtiers par \vp{VETTWIN_RELAY_KEY}.

Lancement au démarrage :
\begin{itemize}
  \item \textbf{Raspberry Pi, Linux} : \vp{salle/vettwins-relay.service} (remplacer \aremplacer{UTILISATEUR}) puis
  \cmd{sudo systemctl enable --now vettwins-relay} ; journal : \cmd{journalctl -u vettwins-relay -f}.
  \item \textbf{Mac mini} : \vp{salle/fr.vettwins.relay.plist} (remplacer \aremplacer{VOUS}), comme le service web
  de \vp{DEPLOIEMENT_MACMINI}. Le relais sert lui-même \vp{web/} : le service \vp{fr.vettwins.web} n'est pas
  nécessaire dans la salle. Si le pare-feu macOS est actif, autoriser les connexions entrantes pour \vp{python3}.
\end{itemize}

\subsection{Contrôle}
\begin{Verbatim}
curl -s http://192.168.50.10:8765/api/status
{"ok": true, "protocol": 1, "regie": false, "clients": {}, "last_snapshot_age_s": null}
\end{Verbatim}
\attention{Ne pas exposer le relais sur Internet (pas de \cmd{tailscale funnel}, pas de redirection de port) : l'API
des boîtiers n'a pas d'authentification forte, seulement une clé partagée en clair.}

\section{Poste de régie}
\begin{enumerate}
  \item Ouvrir \vp{http://192.168.50.10:8765/index.html?room=1} (ajouter \vp{&key=...} si besoin), choisir le patient
  dans la barre latérale (poids, maladie), puis l'onglet \emph{Anesthésie}.
  \item La ligne \emph{Salle} indique : relais connecté, capteurs posés (barrés s'ils sont absents), site du
  stéthoscope, et donne les liens du moniteur de salle et du lecteur d'auscultation.
  \item Choisir le scénario chirurgical, démarrer la session. Laisser la vitesse à \(\times 1\) pendant la séance.
  \item Chaque action d'un boîtier apparaît dans la barre d'état, préfixée par son nom (\cmd{[pupitre] propofol 3 mg/kg IV}).
  La feuille d'anesthésie et le journal la conservent pour le débriefing.
\end{enumerate}
Une seule régie à la fois : si une seconde page \vp{?room=1} se connecte, la première est prévenue et ne reçoit plus
les boîtiers. Les entrées des boîtiers sont refusées tant qu'aucune session n'est démarrée, sauf les capteurs et le
stéthoscope, qui sont mémorisés.

\section{Mise en route d'un boîtier ESP32 (commun)}\label{sec:esp32}
Les six boîtiers utilisent une carte \textbf{ESP32 DevKitC} (module ESP32-WROOM-32, port USB), alimentée par un
chargeur USB 5~V ou une batterie USB. Les entrées sont actives à l'état bas : chaque contact ou bouton relie sa
broche à la masse (GND), la résistance de tirage interne de l'ESP32 suffit.

\begin{enumerate}
  \item Installer l'IDE Arduino 2. \emph{Préférences} \(\to\) URL de cartes supplémentaires :\\
  \vp{https://espressif.github.io/arduino-esp32/package_esp32_index.json} ; \emph{Gestionnaire de cartes} :
  « esp32 » (Espressif). Carte : \emph{ESP32 Dev Module}.
  \item Copier \vp{salle/firmware/libraries/VetTwinRelay} dans le dossier \vp{libraries} du carnet de croquis
  (\vp{~/Documents/Arduino/libraries/} sous macOS et Windows, \vp{~/Arduino/libraries/} sous Linux).
  \item Dans cette copie, dupliquer \vp{src/VetTwinConfig.example.h} en \vp{src/VetTwinConfig.h} et le remplir
  (réseau de la salle, adresse du relais, clé). Ce fichier n'est pas versionné (mot de passe Wi-Fi).
  \item Pour le pupitre : installer \emph{MFRC522} et \emph{LiquidCrystal I2C} (Frank de Brabander) par le
  gestionnaire de bibliothèques.
  \item Ouvrir le croquis (\texttt{salle/firmware/<boîtier>/<boîtier>.ino}), téléverser, ouvrir le moniteur série à
  115\,200 bauds : le boîtier affiche son adresse IP puis chaque envoi et son verdict.
\end{enumerate}
Équivalent en ligne de commande :
\begin{Verbatim}
arduino-cli core install esp32:esp32 \
  --additional-urls https://espressif.github.io/arduino-esp32/package_esp32_index.json
arduino-cli lib install MFRC522 "LiquidCrystal I2C"
arduino-cli compile --fqbn esp32:esp32:esp32 --libraries salle/firmware/libraries salle/firmware/capteurs
arduino-cli upload  --fqbn esp32:esp32:esp32 -p /dev/ttyUSB0 salle/firmware/capteurs
\end{Verbatim}

\textbf{Broches à éviter} sur l'ESP32 : GPIO 6 à 11 (mémoire flash), GPIO 0, 2, 12 et 15 au démarrage si on y relie un
contact fermé au repos (broches de configuration), GPIO 34 à 39 n'ont pas de résistance de tirage (réservées ici
aux potentiomètres). Avec le Wi-Fi actif, seules les entrées analogiques de l'ADC1 (GPIO 32 à 39) fonctionnent.

\textbf{Essai sans boîtier} : chaque entrée se simule avec \vp{curl} depuis n'importe quel ordinateur du réseau,
ce qui permet de préparer une séance avant que le matériel soit prêt :
\begin{Verbatim}
curl -s -X POST http://192.168.50.10:8765/api/device -H 'Content-Type: application/json' \
  -d '{"device":"essai","input":{"type":"probes","spo2":false}}'
{"ok": true, "text": "capteurs : ecg ✓ spo2 ✗ nibp ✓ art ✓ capno ✓ temp ✓"}
\end{Verbatim}
Codes HTTP : 200 (verdict du moteur, champ \vp{ok}), 400 (entrée mal formée), 403 (clé), 503 (régie absente),
504 (régie muette plus de 3~s).

%---------------------------------------------------------------------------------------------
\section{Dispositif 1 : écran moniteur de la salle}
\subsection{Rôle}
Afficher aux étudiants le moniteur d'anesthésie seul (ECG D~II, pression artérielle, capnogramme, valeurs, alarmes,
bip à chaque QRS de tonalité liée à la SpO\(_2\)), sans les commandes ni les variables cachées du modèle (plan,
concentrations, cause d'un décès). Une voie dont le capteur n'est pas posé (dispositif~2) reste vide, et ses alarmes
ne sont pas affichées.

\subsection{Matériel}
\begin{itemize}
  \item Écran de 24 à 32 pouces sur bras articulé ou pied roulant, à la tête du mannequin.
  \item Ordinateur d'affichage : mini-PC, Raspberry Pi 4 ou 5 (le moniteur ne calcule pas, il dessine), ou tablette.
  \item Haut-parleurs actifs (bips et alarmes).
\end{itemize}

\subsection{Mise en service}
Adresse : \vp{http://192.168.50.10:8765/monitor.html?room=1&autostart=1}. Le navigateur exige un clic pour
autoriser le son ; en mode kiosque, Chromium peut l'autoriser d'office :
\begin{Verbatim}
chromium --kiosk --autoplay-policy=no-user-gesture-required \
  "http://192.168.50.10:8765/monitor.html?room=1&autostart=1"
\end{Verbatim}
(\vp{chromium-browser} sur certaines versions de Raspberry Pi OS ; à placer dans le démarrage automatique de la
session graphique, par exemple \vp{~/.config/labwc/autostart} sous Raspberry Pi OS Bookworm.) Désactiver la mise en
veille de l'écran. Sans \vp{autostart=1}, un bouton « Démarrer le moniteur » s'affiche.

En-tête : « en ligne » (vert) si la régie diffuse ; « régie absente » ou « en attente de la régie » sinon. Le
moniteur se reconnecte seul après une coupure du relais ou de la régie, et reçoit à la connexion le dernier état.

\subsection{Limites}
Le capnogramme est schématique (plateau à l'EtCO\(_2\) au rythme respiratoire), comme dans l'onglet de la régie. La PNI
est affichée en continu (pas de cycle de gonflage). Ces limites sont celles de l'onglet Anesthésie
(\vp{docs/ANESTHESIE.tex}).

%---------------------------------------------------------------------------------------------
\section{Dispositif 2 : boîtier des capteurs posés}
\subsection{Rôle}
Savoir quels capteurs l'étudiant a réellement posés sur le mannequin : électrodes ECG, pince de l'oxymètre, brassard,
cathéter artériel, capnographe (raccord sur la sonde), sonde thermique. Entrée \vp{probes}.

\subsection{Matériel}
\begin{itemize}
  \item 1 ESP32 DevKitC, 1 boîtier, câble souple à 2 conducteurs par capteur.
  \item Un contact par capteur, au choix : \textbf{interrupteur à lame souple} (ILS, normalement ouvert) sous la peau
  du mannequin et un petit aimant néodyme dans le capteur factice ; ou \textbf{micro-rupteur} fermé par le capteur
  (brassard serré, raccord vissé).
\end{itemize}

\subsection{Câblage}
\begin{tabularx}{\linewidth}{lllL}
\toprule
\textbf{Capteur} & \textbf{Clé} & \textbf{Broche} & \textbf{Réalisation suggérée} \\
\midrule
électrodes ECG & \vp{ecg} & GPIO 13 & 3 pinces en série (contact fermé quand les 3 sont posées), ou ILS sous le patch du membre de référence \\
oxymètre & \vp{spo2} & GPIO 14 & ILS dans la langue ou le pli du flanc, aimant dans la pince \\
brassard & \vp{nibp} & GPIO 27 & micro-rupteur sous le velcro, ou ILS + aimant dans le brassard \\
cathéter artériel & \vp{art} & GPIO 26 & ILS au point de ponction (artère pédieuse), aimant sur l'embase du cathéter \\
capnographe & \vp{capno} & GPIO 25 & ILS dans le raccord de la sonde, aimant dans l'adaptateur \\
sonde thermique & \vp{temp} & GPIO 33 & ILS dans le canal \og œsophagien \fg{} ou rectal \\
LED témoin & & GPIO 2 & LED de la carte : dernier envoi accepté \\
\bottomrule
\end{tabularx}
L'autre borne de chaque contact va à GND. Contact fermé = capteur posé.

\subsection{Code}
\vp{salle/firmware/capteurs/capteurs.ino} (annexe~\ref{app:code}). À chaque changement, l'état des six capteurs est
envoyé ; il est renvoyé toutes les 30~s (une régie redémarrée retrouve l'état) et 2~s après un échec.

\subsection{Essai}
Poser puis retirer chaque capteur : la ligne \emph{Salle} de la régie barre ou rétablit la clé, la voie correspondante
du moniteur se vide ou se remplit. Sans boîtier des capteurs, tout est affiché.

%---------------------------------------------------------------------------------------------
\section{Dispositif 3 : auscultation}
\subsection{Rôle}
L'étudiant pose un vrai stéthoscope sur le thorax du mannequin et entend les bruits du cœur du patient virtuel, au foyer
où il le pose (base ou apex), à la FC courante de la session d'anesthésie.

\subsection{Matériel}
\begin{itemize}
  \item \textbf{Son} : un transducteur de surface (excitateur, 5 à 10~W) collé à l'intérieur de la paroi thoracique,
  ou un petit haut-parleur large bande de 40 à 60~mm contre la paroi ; un amplificateur de classe D 5~V
  (type PAM8403) ; une carte son USB.
  \item \textbf{Poste de lecture} : l'ordinateur qui joue le son ouvre \vp{ausc.html}. Le plus simple est un
  Raspberry Pi ou un mini-PC dédié, sous la table ou dans le mannequin, dont la sortie audio par défaut est
  l'amplificateur. Ce poste exécute son propre moteur (API \vp{pcg}) : temps de calcul sur Raspberry Pi non mesuré.
  \item \textbf{Détection du foyer} : 1 ESP32, 2 capteurs à effet Hall omnipolaires fonctionnant en 3,3~V à sortie
  collecteur ouvert (par exemple Honeywell SS451A ; l'A3144 demande au moins 4,5~V et ne convient pas), un aimant
  néodyme (\(\varnothing\)10~mm environ) collé au dos du pavillon du stéthoscope d'entraînement.
\end{itemize}

\subsection{Câblage}
\begin{tabularx}{\linewidth}{lllL}
\toprule
\textbf{Élément} & \textbf{Broche} & \textbf{Alimentation} & \textbf{Position} \\
\midrule
Hall « base » & GPIO 32 & 3V3, GND & foyer de la base (côté gauche, haut du cœur) \\
Hall « apex » & GPIO 33 & 3V3, GND & foyer de l'apex (côté gauche, choc de pointe) \\
\bottomrule
\end{tabularx}
Sortie de la carte son USB \(\to\) entrée de l'amplificateur \(\to\) transducteur. Alimenter l'amplificateur par le même
bloc USB que le poste de lecture évite les boucles de masse (ronflement).

\subsection{Code et mise en service}
\begin{itemize}
  \item \vp{salle/firmware/stethoscope/stethoscope.ino} : envoie \vp{{"type":"stethoscope","site":"apex"}}, ou
  \vp{null} 0,6~s après que l'aimant a quitté les deux capteurs (on peut glisser d'un foyer à l'autre sans coupure).
  \item Poste de lecture :
\begin{Verbatim}
chromium --kiosk --autoplay-policy=no-user-gesture-required \
  "http://192.168.50.10:8765/ausc.html?room=1&autostart=1"
\end{Verbatim}
  La page boucle 8~s de son synthétisé et en recalcule une quand la FC change d'au moins 3~/min, avec un fondu de
  0,4~s. Le menu \emph{Sortie} choisit une autre sortie audio (Chromium récent ; sinon, régler la sortie par défaut
  du système). Les boutons \emph{base}, \emph{apex}, \emph{silence} servent à régler le volume sans stéthoscope.
\end{itemize}

\subsection{Limites}
\begin{itemize}
  \item Les agents d'anesthésie ne changent que la \emph{fréquence} du son : souffle et timbre sont ceux de la maladie du
  patient, pas ceux de l'état hémodynamique de la session.
  \item Limites du module d'auscultation (\vp{docs/AUSCULTATION.md}) : calage sur huit enregistrements (un chien par
  grade de souffle), intervalle B1--B2 trop long (0,26~s contre 0,18~s mesuré à FC 90), attribution du site « base »
  à revoir. Domaine du synthétiseur : FC de 20 à 300~/min ; au-dehors, et après un décès, silence.
  \item Le transducteur ne reproduit pas fidèlement les basses fréquences de B1 et B2 à travers la paroi : régler le
  volume à l'oreille, avec le stéthoscope.
\end{itemize}

%---------------------------------------------------------------------------------------------
\section{Dispositif 4 : pupitre d'injection RFID}
\subsection{Rôle}
Remplacer le bouton « Injecter » de l'écran par un geste : l'étudiant prépare une vraie seringue (remplie d'eau),
la présente au pupitre, règle le volume poussé, injecte. Le pupitre affiche le verdict du moteur.

\subsection{Matériel}
\begin{itemize}
  \item 1 ESP32 DevKitC ; 1 lecteur RFID RC522 (13,56~MHz) ; étiquettes NFC autocollantes NTAG213 ou NTAG215
  (une par seringue d'entraînement) ;
  \item 1 bouton rotatif à crans (module KY-040) ; 2 boutons poussoirs (dont un gros bouton « INJECTER ») ;
  \item 1 écran LCD 16\(\times\)2 avec adaptateur I2C (PCF8574, adresse 0x27 ou 0x3F) ;
  \item 1 boîtier ; en option, un réservoir (\og veine \fg{} : flacon et tubulure avec site d'injection) pour pousser
  réellement le liquide.
\end{itemize}

\subsection{Câblage}
\begin{tabularx}{\linewidth}{lllL}
\toprule
\textbf{Module} & \textbf{Broche du module} & \textbf{ESP32} & \textbf{Remarque} \\
\midrule
RC522 & SDA (SS) / SCK / MOSI / MISO & GPIO 5 / 18 / 23 / 19 & bus SPI matériel \\
RC522 & RST / 3.3V / GND & GPIO 4 / 3V3 / GND & \textbf{3,3~V uniquement} \\
LCD I2C & SDA / SCL & GPIO 21 / 22 & voir l'encadré ci-dessous \\
KY-040 & CLK / DT / SW & GPIO 32 / 33 / 25 & SW : pas de 0,1 ou 1~mL \\
bouton INJECTER & & GPIO 26 & vers GND \\
bouton ANNULER & & GPIO 27 & oublie la seringue présentée \\
\bottomrule
\end{tabularx}
\attention{Le RC522 se détruit en 5~V. L'adaptateur I2C du LCD s'alimente le plus souvent en 5~V et tire alors SDA
et SCL à 5~V : intercaler un convertisseur de niveau I2C bidirectionnel (BSS138), ou choisir un écran prévu pour 3,3~V.}

\subsection{Table des seringues}
Les étiquettes ne contiennent rien : le pupitre lit leur identifiant (UID) et cherche la seringue dans la table
\vp{SYRINGES} du croquis (agent, mg/mL, voie, durée d'injection IV, libellé). Pour enregistrer une étiquette :
la présenter, noter l'UID affiché (écran et moniteur série), l'ajouter à la table, téléverser. Les concentrations
proposées sont celles des présentations commerciales de l'onglet Anesthésie (\vp{docs/ANESTHESIE_SOURCES.md}) ;
les identifiants d'agent sont ceux des modules (\vp{modules/anesthesia/agents/}). Les UID de la table livrée sont
des exemples à remplacer.

Voies : la régie refuse une voie absente du module (propofol IM, par exemple). Voies disponibles : IV pour tous ;
IM aussi pour acépromazine, atipamézole, dexmédétomidine, lidocaïne, méthadone, midazolam, naloxone ; SC pour la
méthadone.

\subsection{Utilisation}
Présenter la seringue (l'écran affiche son libellé), tourner le bouton pour le volume (appui sur le bouton : pas de
1~mL), appuyer sur INJECTER. Entrée envoyée :
\begin{Verbatim}
{"type":"inject","agent":"propofol","conc_mg_ml":10,"ml":6,"route":"iv","duration_s":30}
\end{Verbatim}
La régie calcule \(6 \times 10 / \text{poids}\) mg/kg ; l'écran affiche \og OK \fg{} et l'action, ou \og REFUSE \fg{}
et le motif (accents retirés pour l'écran). Une seringue oubliée plus de 2~minutes est désélectionnée.

\subsection{Limites}
Le volume est \emph{déclaré} au bouton, pas mesuré ; le pupitre ne suit pas le volume restant dans la seringue.
Mesurer la course du piston (capteur de distance à temps de vol, potentiomètre linéaire) est une évolution possible.
Une étiquette recollée sur une autre seringue trompe le pupitre : marquer les seringues.

%---------------------------------------------------------------------------------------------
\section{Dispositif 5 : machine d'anesthésie factice}
\subsection{Rôle}
Régler le vaporisateur, le débit de gaz frais et intuber sur un objet, et non plus à l'écran. Entrées \vp{vaporizer},
\vp{fresh_gas}, \vp{intubate}, \vp{extubate}.

\subsection{Matériel}
\begin{itemize}
  \item 1 ESP32 DevKitC ; 2 potentiomètres rotatifs linéaires 10~k\(\Omega\) (axe de 6~mm) ; 1 interrupteur à bascule
  (choix de l'agent) ; 1 ILS et 1 aimant annulaire ; 2 LED (verte, rouge) avec résistances de 330~\(\Omega\).
  \item Façade : cadran de vaporisateur imprimé en 3D ou vaporisateur réformé dont on a retiré la cuve, bouton de
  débitmètre.
\end{itemize}

\subsection{Câblage}
\begin{tabularx}{\linewidth}{llLL}
\toprule
\textbf{Élément} & \textbf{Broche} & \textbf{Branchement} & \textbf{Remarque} \\
\midrule
potentiomètre vaporisateur & GPIO 34 & extrémités sur 3V3 et GND, curseur sur la broche & ADC1 (compatible Wi-Fi) \\
potentiomètre gaz frais & GPIO 35 & idem & ADC1 \\
sélecteur d'agent & GPIO 16 & vers GND & ouvert = isoflurane, fermé = sévoflurane \\
ILS d'intubation & GPIO 17 & vers GND & dans la trachée du mannequin, aimant sur la sonde \\
LED verte / rouge & GPIO 18 / 19 & anode sur la broche, cathode vers GND par 330~\(\Omega\) & verdict \\
\bottomrule
\end{tabularx}

\subsection{Graduation et code}
\vp{salle/firmware/machine/machine.ino}. Les butées du cadran sont les \vp{vaporizer_max_pct} des modules :
5~\% pour l'isoflurane, 8~\% pour le sévoflurane ; graduation de 0,25~\% comme l'onglet Anesthésie ; les 3~\%
du bas de course forment le cran \og fermé \fg{}. Graduer la façade en conséquence (angle utile du potentiomètre
divisé en 20 ou 32 pas). Le débitmètre va de 0 à 4~L/min par 0,25 : c'est une pleine échelle choisie pour le boîtier,
à modifier dans \vp{FGF_MAX}. Un réglage n'est envoyé qu'après 400~ms de stabilité.

Si le cadran dépasse la butée de l'agent choisi, la régie ramène à la butée et le signale. Changer d'agent avec le
vaporisateur ouvert ferme l'ancien.

\subsection{Intubation}
Sonde en place (aimant devant l'ILS) : \vp{intubate}. Le moteur refuse si le patient déglutit encore (« intubation
impossible : déglutition présente (plan trop léger) ») : la LED rouge s'allume, l'étudiant retire la sonde et
recommence plus tard. Sonde retirée : \vp{extubate}, qui arrête aussi la ventilation contrôlée.

\subsection{Limites}
Pas de ballon ni de ventilateur sur le boîtier : la ventilation contrôlée se règle à la régie. Pas de modèle des
voies aériennes (intubation sélective, obstruction) dans le moteur.

%---------------------------------------------------------------------------------------------
\section{Dispositif 6 : pouls fémoral}
\subsection{Rôle}
Boîtier de \emph{sortie} : un vibreur sous la peau de la face interne de la cuisse bat à la FC du patient virtuel, avec
une force qui suit la pression pulsée (PAS \(-\) PAD). Plus de pouls si le patient est décédé, si la régie est absente
ou si l'état a plus de 5~s.

\subsection{Matériel}
\begin{itemize}
  \item 1 ESP32 DevKitC ; 1 moteur vibrant à masselotte de type « pièce » (ERM, 3~V, \(\varnothing\)10~mm), ou un
  mini-solénoïde 5~V pour un choc plus net ;
  \item 1 MOSFET canal N à seuil logique (IRLZ44N en boîtier TO-220, ou AO3400) ; 1 diode de roue libre (1N5819) ;
  résistances de 100~\(\Omega\) (grille) et 100~k\(\Omega\) (grille vers GND).
\end{itemize}

\subsection{Câblage}
\begin{Verbatim}
+3V3 (ou +5V pour un solénoïde) ----+---- moteur (+)
                                    |
                           cathode 1N5819 (bague)
                           anode   1N5819
                                    |
                                    +---- moteur (-) ---- drain du MOSFET
GPIO 25 ---- 100 ohms ---- grille du MOSFET
grille du MOSFET ---- 100 kohms ---- GND
source du MOSFET ---- GND
\end{Verbatim}
La diode, montée en sens inverse aux bornes du moteur, protège le MOSFET des surtensions. Un moteur de 3~V se
branche sur la sortie 3V3 de la carte (quelques dizaines de mA) ; un solénoïde sur le 5~V USB.

\subsection{Code et réglage}
\vp{salle/firmware/pouls/pouls.ino}. Une tâche lit \cmd{GET /api/pulse} chaque seconde (ligne
\texttt{fc pas pad décès valide âge}) ; la boucle principale produit les battements à la FC. Deux réglages d'usage, à
faire au toucher : \vp{PP_FULL} (pression pulsée qui donne la force maximale, 60~mmHg par défaut) et
\vp{MOTOR_START} (rapport cyclique en dessous duquel ce moteur ne démarre pas). Ce ne sont pas des seuils
cliniques : aucun seuil de \og pouls non palpable \fg{} n'est sourcé dans le projet.

\subsection{Limites}
Les battements sont produits localement à la FC : ils ne sont pas synchrones des QRS du moniteur, et une arythmie
(bloc AV) n'apparaît pas dans le pouls. La FC est mise à jour chaque seconde.

%---------------------------------------------------------------------------------------------
\section{Dispositif 7 : boîtier du formateur}
\subsection{Rôle}
Déclencher les temps du scénario depuis la salle, sans retourner à la régie : temps opératoire suivant, hémorragie,
stimulus (pince), bolus de cristalloïde, fin de chirurgie.

\subsection{Matériel et câblage}
1 ESP32 DevKitC, 5 boutons d'arcade (30~mm), 2 LED avec résistances de 330~\(\Omega\), 1 boîtier.

\begin{tabularx}{\linewidth}{llL}
\toprule
\textbf{Bouton} & \textbf{Broche} & \textbf{Entrée envoyée} \\
\midrule
temps opératoire suivant & GPIO 13 & \vp{{"type":"procedure_next"}} (le suivant du scénario choisi à la régie) \\
hémorragie & GPIO 14 & \vp{{"type":"hemorrhage","ml_kg":5}} \\
stimulus & GPIO 27 & \vp{{"type":"stimulus","intensity":1,"duration_s":30}} \\
bolus de cristalloïde & GPIO 26 & \vp{{"type":"fluid_bolus","ml_kg":10}} \\
fin de chirurgie & GPIO 25 & \vp{{"type":"end_surgery"}} \\
LED verte / rouge & GPIO 18 / 19 & verdict \\
\bottomrule
\end{tabularx}
Les valeurs sont celles proposées par défaut dans l'onglet Anesthésie ; ce sont des réglages du scénario, à adapter
dans \vp{salle/firmware/formateur/formateur.ino}. Le compteur des temps opératoires suit aussi les boutons de la
régie ; à la fin du scénario, la LED rouge s'allume (« tous les temps opératoires sont réalisés »).

%---------------------------------------------------------------------------------------------
\section{Essais et vérifications}\label{sec:essais}
\subsection{Sur un ordinateur (avant tout montage)}
\begin{Verbatim}
python3 tests/check_relay.py         # relais : 39 contrôles
python3 tests/check_firmware.py      # six croquis compilés sur l'hôte + bibliothèque VetTwinRelay
sh web/build.sh && node tests/room_e2e.mjs   # chaîne complète dans Chromium (Playwright)
\end{Verbatim}

\subsection{Première mise en route de chaque boîtier}
\begin{enumerate}
  \item Moniteur série à 115\,200 bauds : \og Wi-Fi connecté, IP ... \fg{}. Sinon, vérifier le réseau 2,4~GHz et
  \vp{VetTwinConfig.h}.
  \item Régie fermée : une action du boîtier doit donner \og 503 régie non connectée \fg{} dans le moniteur série
  (le relais est joint, la régie non).
  \item Régie ouverte, session démarrée : l'action apparaît dans la barre d'état de la régie, préfixée par le nom du
  boîtier, et le verdict revient au boîtier (LED, écran, moniteur série).
\end{enumerate}

\subsection{Avant chaque séance}
\begin{itemize}
  \item Relais : \cmd{curl http://IP:8765/api/status} ; régie : ligne \emph{Salle} \og relais connecté \fg{}.
  \item Moniteur et lecteur d'auscultation \og en ligne \fg{} ; volume réglé avec les boutons \emph{base} et \emph{apex}.
  \item Capteurs retirés : voies vides au moniteur. Seringues présentées une à une : libellés corrects.
  \item Session démarrée à \(\times 1\), patient et scénario choisis ; feuille d'anesthésie vide.
\end{itemize}

\subsection{Dépannage}
\begin{tabularx}{\linewidth}{>{\raggedright\arraybackslash}p{0.36\linewidth}L}
\toprule
\textbf{Symptôme} & \textbf{Cause probable et remède} \\
\midrule
Boîtier : \og Wi-Fi absent \fg{} & réseau 5~GHz seulement, ou réseau à identifiant (WPA2-Entreprise) : point d'accès dédié en 2,4~GHz. \\
Boîtier : \og relais injoignable \fg{} & adresse ou port faux dans \vp{VetTwinConfig.h} ; relais lancé avec \cmd{--host 127.0.0.1} ; pare-feu du serveur. \\
503 \og régie non connectée \fg{} & page de régie fermée, ou ouverte sans \vp{?room=1}. \\
504 \og pas de réponse de la régie \fg{} & onglet de la régie en arrière-plan sur un portable en veille ; garder la régie au premier plan, désactiver la veille. \\
403 & clé du relais différente entre le relais, les pages (\vp{&key=}) et \vp{VetTwinConfig.h}. \\
\og session non démarrée \fg{} & démarrer la session dans l'onglet Anesthésie ; changer le patient de la barre latérale arrête la session. \\
Moniteur muet & clic \og Démarrer le moniteur \fg{} non fait, ou Chromium sans \vp{--autoplay-policy=no-user-gesture-required}. \\
Pupitre : écran vide & adresse I2C 0x3F au lieu de 0x27 ; contraste (potentiomètre de l'adaptateur). \\
Pupitre : aucune lecture RFID & RC522 alimenté en 5~V (détruit) ou fil SDA/RST inversé. \\
Pouls permanent ou absent & régie absente (pas de pouls, voulu) ; \vp{MOTOR_START} trop haut pour le moteur. \\
\bottomrule
\end{tabularx}

\section{Sécurité et limites générales}
\begin{itemize}
  \item Tout fonctionne en très basse tension (USB 5~V, 3,3~V) ; pas de secteur dans le mannequin. Fixer les câbles
  pour qu'on ne puisse pas tirer sur une carte ; isoler les soudures (gaine thermorétractable).
  \item Les liquides (seringues remplies d'eau, réservoir) restent hors des boîtiers électroniques.
  \item Les valeurs du modèle sont pédagogiques et, pour beaucoup, des placeholders (niveau D) ; le dire aux étudiants
  pendant le briefing. Limites du module d'anesthésie : \vp{docs/ANESTHESIE.tex} (pas de chat, pas de voies
  aériennes, pas de toxicité de la lidocaïne, spirale de décompensation possible après un propofol très rapide).
  \item Un essai reproduit la même séance avec la même graine et les mêmes gestes ; la feuille d'anesthésie et le
  journal de la régie (onglet Anesthésie, \og Actualiser \fg{}) servent au débriefing.
\end{itemize}

\appendix
\section{Protocole des boîtiers (version 1)}
\cmd{POST /api/device} avec \cmd{{"device": "<nom>", "input": {"type": ..., ...}}}. Réponse :
\cmd{{"ok": true|false, "text": "..."}}.

\begin{tabularx}{\linewidth}{>{\raggedright\arraybackslash}p{0.2\linewidth}>{\raggedright\arraybackslash}p{0.36\linewidth}L}
\toprule
\textbf{\vp{type}} & \textbf{Champs} & \textbf{Traduction par la régie} \\
\midrule
\vp{probes} & \vp{ecg}, \vp{spo2}, \vp{nibp}, \vp{art}, \vp{capno}, \vp{temp} (booléens, facultatifs) & capteurs posés (hors moteur) \\
\vp{stethoscope} & \vp{site} : \vp{"base"}, \vp{"apex"} ou \vp{null} & site joué par le lecteur (hors moteur) \\
\vp{inject} & \vp{agent}, \vp{conc_mg_ml}, \vp{ml}, \vp{route} (\vp{iv}), \vp{duration_s} (30) & \vp{bolus} en mg/kg = mL \(\times\) mg/mL / poids \\
\vp{vaporizer} & \vp{agent} (celui de la machine), \vp{pct} & \vp{vaporizer}, borné à \vp{vaporizer_max_pct} \\
\vp{fresh_gas} & \vp{l_min} & \vp{fresh_gas} \\
\vp{oxygen} & \vp{on} & \vp{oxygen} (masque) \\
\vp{intubate}, \vp{extubate} & & idem \\
\vp{procedure_next} & & \vp{procedure_step} suivant du scénario \\
\vp{end_surgery} & & idem \\
\vp{stimulus} & \vp{intensity} (1), \vp{duration_s} (30), \vp{territory} & idem \\
\vp{hemorrhage}, \vp{fluid_bolus} & \vp{ml_kg} & idem \\
\vp{fluids} & \vp{ml_kg_h} & idem \\
\vp{warming} & \vp{on} & idem \\
\vp{circuit} & \vp{kind} : \vp{circle}, \vp{non_rebreathing} & idem \\
\vp{ventilation} & \vp{mode} : \vp{spontaneous} ou \vp{controlled} avec \vp{rr}, \vp{vt_ml_kg} & idem \\
\bottomrule
\end{tabularx}

Autres points de l'API : \cmd{GET /api/status} (régie, clients), \cmd{GET /api/state} (dernier état : pouls, capteurs,
temps), \cmd{GET /api/pulse} (ligne \texttt{fc pas pad décès valide âge\_s}). WebSocket \vp{/ws} : premier message
\vp{{"kind":"hello","role":"regie"|"monitor"|"ausc"}} ; la régie diffuse \vp{snapshot}, \vp{waves}, \vp{session},
\vp{probes}, \vp{steth} et répond aux \vp{device} par \vp{result}.

\section{Code des boîtiers}\label{app:code}
\subsection{Bibliothèque commune \texttt{VetTwinRelay}}
\VerbatimInput{../salle/firmware/libraries/VetTwinRelay/src/VetTwinRelay.h}
\VerbatimInput{../salle/firmware/libraries/VetTwinRelay/src/VetTwinConfig.example.h}
\subsection{Boîtier des capteurs}
\VerbatimInput{../salle/firmware/capteurs/capteurs.ino}
\subsection{Stéthoscope}
\VerbatimInput{../salle/firmware/stethoscope/stethoscope.ino}
\subsection{Pupitre d'injection}
\VerbatimInput{../salle/firmware/pupitre/pupitre.ino}
\subsection{Machine d'anesthésie}
\VerbatimInput{../salle/firmware/machine/machine.ino}
\subsection{Pouls fémoral}
\VerbatimInput{../salle/firmware/pouls/pouls.ino}
\subsection{Boîtier du formateur}
\VerbatimInput{../salle/firmware/formateur/formateur.ino}

\end{document}
